\documentclass[12pt,a4paper]{article}

% ===== Pacotes =====
\usepackage[utf8]{inputenc}
\usepackage[T1]{fontenc}
\usepackage[brazil]{babel}
\usepackage[margin=2.5cm]{geometry}
\usepackage{lmodern}
\usepackage{graphicx}
\usepackage{indentfirst}
\usepackage{booktabs}
\usepackage{url}
\usepackage{hyperref}
\hypersetup{hidelinks}

% Mostra a imagem se existir; senao mostra uma caixa reservada (nao quebra a compilacao)
\newcommand{\figcond}[2]{%
  \IfFileExists{#1}{\includegraphics[width=\linewidth]{#1}}%
  {\fbox{\parbox[c][4.5cm][c]{0.98\linewidth}{\centering \textit{#2}}}}}

\begin{document}

% ===== Cabecalho =====
\begin{center}
{\large \textbf{Sistema Inteligente de Monitoramento da Qualidade do Ar em \\ Ambientes Internos utilizando IoT}}\\[0.6cm]
Gustavo Giacometti\textsuperscript{1}, Giovanni Pontes\textsuperscript{1}, Wallace Santana\textsuperscript{1}\\[0.2cm]
Orientador: Prof. Wallace Rodrigues de Santana\textsuperscript{1}\\[0.3cm]
\textsuperscript{1}Faculdade de Computação e Informática\\
Universidade Presbiteriana Mackenzie (UPM)\\
Rua da Consolação, 930 -- Consolação, São Paulo -- SP, 01302-907 -- Brasil\\[0.2cm]
\small\texttt{\{10401200, 10402204, [RA\_WALLACE]\}@mackenzista.br}\\
\small\texttt{1165744@mackenzie.br}
\end{center}
\vspace{0.4cm}

% ===== Resumos =====
\noindent\textbf{Abstract.} \textit{This paper presents an intelligent indoor air quality monitoring
system based on the Internet of Things (IoT), composed of two simulated devices: one with sensors
(MQ-135 and DHT22) and another with an actuator. An ESP32 reads CO\textsubscript{2}, temperature and
humidity and publishes them as JSON through the MQTT protocol using the Eclipse Mosquitto broker.
Node-RED applies business rules that classify the environment as good, moderate or critical, command
the actuator (ventilation/alert), consume the external WAQI air quality API and persist the readings
into a cloud time-series database (InfluxDB Cloud). Management dashboards are built in Grafana. Results
obtained in a simulated environment (Wokwi) show that the system reliably detects environmental
variations, stores the data and reacts through the actuator, contributing to UN Sustainable
Development Goals 3 and 11.}

\vspace{0.3cm}
\noindent\textbf{Resumo.} Este artigo apresenta um sistema inteligente de monitoramento da qualidade do
ar em ambientes internos baseado em Internet das Coisas (IoT), composto por dois dispositivos simulados:
um com sensores (MQ-135 e DHT22) e outro com um atuador. Um ESP32 lê CO\textsubscript{2}, temperatura e
umidade e os publica em formato JSON por meio do protocolo MQTT utilizando o broker Eclipse Mosquitto. O
Node-RED aplica regras de negócio que classificam o ambiente como bom, moderado ou crítico, comandam o
atuador (ventilação/alerta), consomem a API externa de qualidade do ar WAQI e persistem as leituras em um
banco de dados de séries temporais em nuvem (InfluxDB Cloud). Os \textit{dashboards} gerenciais são
construídos no Grafana. Os resultados obtidos em ambiente simulado (Wokwi) demonstram que o sistema
detecta de forma confiável as variações ambientais, armazena os dados e reage por meio do atuador,
contribuindo para os Objetivos de Desenvolvimento Sustentável 3 e 11 da ONU.

\vspace{0.5cm}

% ===== 1. Introducao =====
\section{Introdução}

A qualidade do ar em \textbf{ambientes internos} --- entendidos como espaços fechados de longa
permanência, como residências, salas de aula, escritórios, hospitais e meios de transporte --- exerce
impacto direto na saúde, no conforto e na produtividade das pessoas. Estima-se que os indivíduos passem
mais de 80\% do tempo em locais fechados, nos quais a concentração de poluentes pode ser de duas a cinco
vezes superior à do ambiente externo \cite{who2021}. Concentrações de dióxido de carbono
(CO\textsubscript{2}) acima de 1000~ppm, decorrentes de ventilação inadequada e alta ocupação, estão
associadas à fadiga, à queda de desempenho cognitivo e ao agravamento de problemas respiratórios
\cite{ashrae2019}. Apesar disso, o monitoramento contínuo dessas condições ainda é pouco difundido.

Com o avanço da \textbf{Internet das Coisas} (IoT), tornou-se viável integrar sensores de baixo custo,
dispositivos embarcados e plataformas de processamento em rede, permitindo a coleta, o armazenamento e a
análise contínuos de dados ambientais, bem como a atuação automática sobre o ambiente \cite{atzori2010}.

\subsection{Descrição e aplicabilidade da solução}

Este trabalho propõe um \textbf{sistema inteligente de monitoramento da qualidade do ar em ambientes
internos} composto por dois dispositivos: um de \textbf{sensoriamento} (CO\textsubscript{2}, temperatura
e umidade) e um \textbf{atuador} que aciona ventilação/alerta quando o ar se torna crítico. O sistema
transmite os dados via MQTT, classifica automaticamente a condição do ambiente, armazena o histórico em
nuvem e apresenta as informações em \textit{dashboards}. A solução é diretamente aplicável a contextos
reais: em \textbf{escolas}, ajuda a manter salas ventiladas e favoráveis ao aprendizado; em
\textbf{escritórios}, contribui para o bem-estar e a produtividade; em \textbf{hospitais}, apoia o
controle de ambientes sensíveis; e em \textbf{residências}, oferece um indicador acessível da salubridade
do lar. Por usar componentes de baixo custo e ferramentas de código aberto, é replicável e escalável.

\subsection{Objetivos}

O objetivo geral é desenvolver uma solução IoT capaz de monitorar variáveis ambientais, classificar a
qualidade do ar, atuar de forma preventiva e disponibilizar as informações em tempo real. Objetivos
específicos: (i) implementar a aquisição com MQ-135 e DHT22; (ii) acionar um atuador via MQTT conforme
regras de negócio; (iii) estabelecer a comunicação via broker Mosquitto; (iv) integrar a API externa
WAQI; (v) persistir os dados em banco de séries temporais em nuvem (InfluxDB) e apresentá-los em
\textit{dashboards} no Grafana.

% ===== 2. ODS =====
\section{Aderência aos Objetivos de Desenvolvimento Sustentável}

O projeto alinha-se à Agenda 2030 da Organização das Nações Unidas \cite{onu2015}, contribuindo para:

\begin{itemize}
  \item \textbf{ODS 3 -- Saúde e Bem-Estar:} ao monitorar a qualidade do ar e atuar sobre condições
  inadequadas, o sistema atua preventivamente contra doenças respiratórias e problemas associados à
  exposição a poluentes, em consonância com a meta 3.9 (reduzir mortes e doenças por poluição do ar).
  \item \textbf{ODS 11 -- Cidades e Comunidades Sustentáveis:} promove ambientes mais saudáveis e
  eficientes, apoiando a meta 11.6 (reduzir o impacto ambiental das cidades). A integração com a API WAQI
  reforça esse vínculo ao comparar o ambiente interno com a qualidade do ar externa da cidade.
\end{itemize}

% ===== 3. Revisao da Literatura =====
\section{Revisão da Literatura}

O conceito de Internet das Coisas, sistematizado por Atzori, Iera e Morabito \cite{atzori2010}, descreve
a interconexão de objetos físicos dotados de sensores e atuadores, formando a base para aplicações de
monitoramento ambiental, tipicamente organizadas em camadas de aquisição, comunicação e processamento.

Quanto aos efeitos da qualidade do ar sobre a saúde, a Organização Mundial da Saúde \cite{who2021}
consolidou diretrizes que reforçam a relevância do monitoramento de poluentes em ambientes internos.
Complementarmente, a norma ANSI/ASHRAE 62.1 \cite{ashrae2019} estabelece parâmetros de ventilação e
referencia o valor de 1000~ppm de CO\textsubscript{2} como indicador de ventilação adequada,
fundamentando os limites adotados neste trabalho.

Na comunicação, o protocolo MQTT \cite{mqtt2014} consolidou-se em IoT por ser leve e baseado no modelo
publicação/assinatura. Na implementação, empregam-se ferramentas de código aberto amplamente adotadas: o
broker Eclipse Mosquitto \cite{mosquitto}, a plataforma de programação \textit{low-code} Node-RED
\cite{nodered}, o banco de séries temporais InfluxDB \cite{influxdata} e a ferramenta de
\textit{dashboards} Grafana \cite{grafana}. Em comparação a soluções comerciais de monitoramento de
qualidade do ar, geralmente de custo elevado e baixa flexibilidade, a abordagem proposta prioriza baixo
custo, abertura e personalização.

% ===== 4. Materiais e Metodos =====
\section{Materiais e Métodos}

O sistema adota uma arquitetura distribuída de IoT, organizada em aquisição/atuação, comunicação,
processamento, armazenamento e visualização, conforme o fluxo:
\begin{center}
Sensores/Atuador (ESP32) $\rightarrow$ MQTT (Mosquitto) $\rightarrow$ Node-RED $\rightarrow$
InfluxDB Cloud $\rightarrow$ Grafana
\end{center}

\begin{figure}[h]
  \centering
  \figcond{arquitetura.png}{[Figura 1 -- Inserir o diagrama de blocos da arquitetura]}
  \caption{Arquitetura do sistema de monitoramento da qualidade do ar.}
  \label{fig:arq}
\end{figure}

\subsection{Dispositivo 1 -- Sensores}
Um microcontrolador \textbf{ESP32} realiza a leitura periódica (a cada 5~s) e formata os dados em JSON. O
\textbf{sensor MQ-135} detecta gases como o CO\textsubscript{2}; o \textbf{sensor DHT22} fornece
temperatura e umidade relativa. No ambiente simulado (Wokwi), a saída analógica do MQ-135 é representada
por um potenciômetro. A programação foi feita em C/C++. O diagrama de conexão é apresentado na
Figura~\ref{fig:dev1}.

\begin{figure}[h]
  \centering
  \figcond{dispositivo1.png}{[Figura 2 -- Inserir o diagrama de conexao do dispositivo 1 (Wokwi)]}
  \caption{Diagrama de conexão do dispositivo de sensoriamento (ESP32 + MQ-135 + DHT22).}
  \label{fig:dev1}
\end{figure}

\subsection{Dispositivo 2 -- Atuador}
Um segundo ESP32 \textbf{assina} um tópico MQTT de comando e aciona um \textbf{atuador} (LED indicador e
\textit{buzzer}), representando um sistema de ventilação/alerta. Quando o Node-RED classifica o ambiente
como crítico, publica o comando \texttt{LIGAR}; caso contrário, \texttt{DESLIGAR}. Assim, o estado do
atuador é \textbf{resultado da leitura de um dado via broker MQTT} a partir de uma regra de negócio.

\begin{figure}[h]
  \centering
  \figcond{dispositivo2.png}{[Figura 3 -- Inserir o circuito do atuador (Wokwi)]}
  \caption{Diagrama de conexão do atuador (ESP32 + LED + \textit{buzzer}).}
  \label{fig:dev2}
\end{figure}

\subsection{Comunicação via MQTT}
Os dispositivos conectam-se à Internet por Wi-Fi e utilizam o broker público \textbf{Eclipse Mosquitto}
(\url{test.mosquitto.org}) \cite{mosquitto}. O dispositivo de sensores publica no tópico
\texttt{mackenzie/oic/qualidadear} e o atuador assina o tópico \texttt{mackenzie/oic/comando}.

\subsection{Processamento, regras de negócio e API}
No \textbf{Node-RED} \cite{nodered}, as leituras são classificadas segundo os limites da
Tabela~\ref{tab:regras}. O Node-RED também consome a API \textbf{WAQI} \cite{waqi}, obtendo a qualidade do
ar externa da cidade de São Paulo como referência comparativa. A Figura~\ref{fig:nodered} apresenta o
fluxo implementado.

\begin{figure}[h]
  \centering
  \figcond{node.png}{[Figura 4 -- Inserir o fluxo do Node-RED]}
  \caption{Fluxo no Node-RED: regras de negócio, comando ao atuador, consumo da API e gravação no InfluxDB.}
  \label{fig:nodered}
\end{figure}

\begin{table}[h]
  \centering
  \caption{Regras de negócio para classificação do ambiente.}
  \label{tab:regras}
  \begin{tabular}{lccc}
    \toprule
    \textbf{Variável} & \textbf{Bom} & \textbf{Moderado} & \textbf{Crítico} \\
    \midrule
    CO\textsubscript{2} (ppm) & $<$ 1000 & 1000 -- 2000 & $>$ 2000 \\
    Temperatura (°C) & \multicolumn{3}{c}{conforto 18 -- 26 (fora da faixa gera alerta)} \\
    Umidade (\%) & \multicolumn{3}{c}{conforto 30 -- 60 (fora da faixa gera alerta)} \\
    \bottomrule
  \end{tabular}
\end{table}

\subsection{Banco de dados em nuvem (InfluxDB Cloud)}
As leituras são gravadas em um banco de \textbf{séries temporais em nuvem}, o \textbf{InfluxDB Cloud}
\cite{influxdata}, no \textit{bucket} \texttt{qualidade\_ar} (\textit{measurement} \texttt{leituras}),
registrando CO\textsubscript{2}, temperatura, umidade, \textit{status} e horário, permitindo o
acompanhamento histórico.

\subsection{Visualização (Grafana)}
Os \textit{dashboards} gerenciais são construídos no \textbf{Grafana} \cite{grafana}, ferramenta
\textbf{externa} conectada ao InfluxDB como fonte de dados (linguagem SQL). Foram criados painéis de
medidor (\textit{gauge}) e de série temporal para CO\textsubscript{2}, temperatura e umidade, além do
indicador de \textit{status} do ambiente.

\subsection{Topologia lógica}
Conforme exigido, cada elemento da aplicação está hospedado em uma máquina distinta. A
Tabela~\ref{tab:topo} apresenta a topologia lógica com os respectivos FQDN e IP.

\begin{table}[h]
  \centering
  \caption{Topologia lógica da aplicação (FQDN e IP).}
  \label{tab:topo}
  \begin{tabular}{lll}
    \toprule
    \textbf{Elemento} & \textbf{FQDN} & \textbf{IP} \\
    \midrule
    Broker MQTT & test.mosquitto.org & 54.36.178.49 \\
    Node-RED & ac07751b50a4 (rede local) & 192.168.0.4 \\
    Banco (InfluxDB Cloud) & us-east-1-1.aws.cloud2.influxdata.com & 34.196.233.7 \\
    API WAQI & api.waqi.info & 139.162.71.178 \\
    Dashboard (Grafana) & [SEU\_USUARIO].grafana.net & (Grafana Cloud) \\
    Dispositivos ESP32 & wokwi.com (simulado) & --- \\
    \bottomrule
  \end{tabular}
\end{table}

\subsection{Implementação e simulação}
Os circuitos foram implementados e validados na plataforma \textbf{Wokwi}, com dois dispositivos: o de
sensoriamento (ESP32 + DHT22 + potenciômetro simulando o MQ-135) e o atuador (ESP32 + LED + \textit{buzzer}).

\vspace{0.2cm}
\noindent\textbf{Links:} repositório: \url{https://github.com/[SEU_USUARIO]/[REPOSITORIO]} \quad
vídeo: \url{https://youtu.be/[ID_DO_VIDEO]}

% ===== 5. Resultados =====
\section{Resultados}

Durante os testes no ambiente simulado, o dispositivo de sensores publicou leituras a cada 5~s, recebidas
e classificadas em tempo real pelo Node-RED e persistidas no InfluxDB Cloud. A Tabela~\ref{tab:cenarios}
resume os cenários avaliados e a classificação obtida.

\begin{table}[h]
  \centering
  \caption{Cenários de teste e classificação obtida.}
  \label{tab:cenarios}
  \begin{tabular}{lcccc}
    \toprule
    \textbf{Cenário} & \textbf{CO\textsubscript{2} (ppm)} & \textbf{Temp. (°C)} & \textbf{Umid. (\%)} & \textbf{Status} \\
    \midrule
    Ambiente ventilado        & 620  & 23,1 & 48 & Bom \\
    Ocupação moderada         & 1480 & 24,3 & 52 & Moderado \\
    Sala fechada e lotada     & 2630 & 27,8 & 64 & Crítico \\
    Recuperação após ventilar & 780  & 23,5 & 50 & Bom \\
    \bottomrule
  \end{tabular}
\end{table}

Em cenários com aumento da concentração de gases, o sistema classificou corretamente o ambiente como
crítico, publicou o comando ao \textbf{atuador} (acionando o LED e o \textit{buzzer} de ventilação/alerta)
e registrou a leitura no InfluxDB. Ao retornar a níveis adequados, o atuador foi desligado
automaticamente, demonstrando a atuação preventiva em malha completa (sensor $\rightarrow$ regra
$\rightarrow$ atuador).

A Figura~\ref{fig:dash} apresenta o \textit{dashboard} construído no Grafana, com os medidores e o
histórico de CO\textsubscript{2}, temperatura e umidade. Sob o ponto de vista gerencial, os painéis
permitem identificar rapidamente o estado do ambiente (pelas faixas de cor) e analisar tendências ao
longo do tempo, apoiando a decisão de ventilar o local.

\begin{figure}[h]
  \centering
  \figcond{dashboard.png}{[Figura 3 -- Inserir o print do dashboard do Grafana]}
  \caption{\textit{Dashboard} de monitoramento no Grafana.}
  \label{fig:dash}
\end{figure}

Os testes evidenciam que a arquitetura proposta é eficiente, escalável e adequada para aplicações reais de
monitoramento ambiental, integrando sensoriamento, atuação, armazenamento em nuvem e visualização externa.

% ===== 6. Conclusoes =====
\section{Conclusões}

O sistema mostrou-se uma solução viável para o monitoramento da qualidade do ar em ambientes internos,
utilizando tecnologias de baixo custo e de código aberto. Os resultados indicam que ele coleta, transmite,
classifica, armazena (em nuvem) e visualiza os dados com eficiência, além de \textbf{atuar} de forma
preventiva por meio do atuador, com contribuição direta para os ODS 3 e 11.

Entre as vantagens destacam-se o baixo custo, a escalabilidade e a flexibilidade da arquitetura. Como
limitações, citam-se a precisão dos sensores de baixo custo e a dependência de conectividade. Como
trabalhos futuros, sugere-se a validação com hardware físico, o uso de sensores calibrados, alertas por
mensageria (e.g., Telegram) e algoritmos de análise preditiva.

% ===== 7. Referencias (ABNT, ordem alfabetica) =====
\begin{thebibliography}{9}

\bibitem{ashrae2019}
AMERICAN SOCIETY OF HEATING, REFRIGERATING AND AIR-CONDITIONING ENGINEERS.
\textbf{ANSI/ASHRAE Standard 62.1}: Ventilation for Acceptable Indoor Air Quality. Atlanta: ASHRAE, 2019.

\bibitem{arduino}
ARDUINO. \textbf{Arduino Documentation}. Disponível em: \url{https://www.arduino.cc}. Acesso em: 2 jun. 2026.

\bibitem{atzori2010}
ATZORI, L.; IERA, A.; MORABITO, G. \textbf{The Internet of Things: a survey}. Computer Networks, v. 54,
n. 15, p. 2787-2805, 2010.

\bibitem{mqtt2014}
BANKS, A.; GUPTA, R. \textbf{MQTT Version 3.1.1}: OASIS Standard. [S.l.]: OASIS, 2014. Disponível em:
\url{https://docs.oasis-open.org/mqtt/mqtt/v3.1.1/os/mqtt-v3.1.1-os.html}. Acesso em: 2 jun. 2026.

\bibitem{mosquitto}
ECLIPSE FOUNDATION. \textbf{Eclipse Mosquitto}: an open source MQTT broker. Disponível em:
\url{https://mosquitto.org}. Acesso em: 2 jun. 2026.

\bibitem{grafana}
GRAFANA LABS. \textbf{Grafana Documentation}. Disponível em: \url{https://grafana.com/docs/}.
Acesso em: 2 jun. 2026.

\bibitem{influxdata}
INFLUXDATA. \textbf{InfluxDB Cloud Documentation}. Disponível em: \url{https://docs.influxdata.com}.
Acesso em: 2 jun. 2026.

\bibitem{nodered}
OPENJS FOUNDATION. \textbf{Node-RED}: low-code programming for event-driven applications. Disponível em:
\url{https://nodered.org}. Acesso em: 2 jun. 2026.

\bibitem{onu2015}
ORGANIZAÇÃO DAS NAÇÕES UNIDAS. \textbf{Transformando Nosso Mundo}: a Agenda 2030 para o Desenvolvimento
Sustentável. Nova York: ONU, 2015. Disponível em: \url{https://brasil.un.org/pt-br/sdgs}.
Acesso em: 2 jun. 2026.

\bibitem{waqi}
WORLD AIR QUALITY INDEX PROJECT. \textbf{Air Quality Open Data Platform}. Disponível em:
\url{https://aqicn.org/api}. Acesso em: 2 jun. 2026.

\bibitem{who2021}
WORLD HEALTH ORGANIZATION. \textbf{WHO global air quality guidelines}. Geneva: WHO, 2021.

\end{thebibliography}

\end{document}
